\documentclass[11pt]{article}
\usepackage[a4paper,margin=0.8in]{geometry}
\usepackage{parskip}
\usepackage{hyperref}
\hypersetup{hidelinks}
\pagestyle{empty}

\begin{document}

\noindent
Quang-Vinh Dang\\
British University Vietnam\\
Hung Yen, Vietnam\\
\href{mailto:vinh.dq4@buv.edu.vn}{vinh.dq4@buv.edu.vn}

\vspace{1em}
\noindent 22 July 2026

\vspace{1em}
\noindent
To the Guest Editors of the Special Issue on\\
\emph{Safety, Alignment, and Responsibility of Large Language Models}\\
(Prof.\ Shuo Shang, Prof.\ Bin Yang, Prof.\ Christian S.\ Jensen, Prof.\ Panos Kalnis)\\
IEEE Transactions on Dependable and Secure Computing

\vspace{1.2em}
\noindent Dear Guest Editors,

\vspace{0.6em}
I am pleased to submit my manuscript, \textbf{``SENTRY: Anytime-Valid E-Process
Monitoring of LLM Agent Action Streams,''} for consideration in the Special
Issue on Safety, Alignment, and Responsibility of Large Language Models.

Large language model agents that call external tools are increasingly deployed
with authority to move money, send messages, and act on a user's behalf, yet
the prevailing runtime defenses threshold heuristic scores at a hand-picked
cut-off and offer \emph{no} statistical guarantee on how often they fire on
benign behavior. This paper addresses that gap directly. SENTRY casts runtime
agent monitoring as sequential change detection and builds an e-process
monitored by an e-detector over per-action surprise signals, controlling the
average run length to false alarm \emph{with no independence assumption} on the
highly dependent action stream---and correcting the learned reference with a
distribution-free PAC threshold so the guarantee holds with finite trusted
data.

The contribution is both methodological and empirical. On real trajectories
collected from a capable agent (DeepSeek-V4-Flash) on the AgentDojo and
$\tau$-bench benchmarks, SENTRY detects $99\%$ of prompt-injection attempts
whose payload is observable in the black-box trace ($128/198$) at a $2\%$
per-step false-alarm rate with a one-action median detection delay, using only
three cheap, model-agnostic signals and no access to model internals. The study
also reframes the guardrail objective from hijack-\emph{success} prediction to
injection-\emph{attempt} detection, and reports a full ablation. A
generalisation study shows the method transfers to a second, unrelated agent
(GPT-4o-mini, AUROC $0.82$) but that its trace-only signal is attack-specific,
failing on InjecAgent's data-integration payloads (AUROC $0.47$)---a boundary I
report and explain rather than hide. To my knowledge this is the first
anytime-valid change-detection study on LLM-agent-safety benchmarks.

The work fits squarely within the Special Issue's scope, in particular
\emph{safety mechanisms in autonomous systems} and \emph{adversarial robustness
and attacks}: it provides a statistically controlled, model-agnostic runtime
tripwire for autonomous LLM agents under indirect prompt injection, together
with a reproducible evaluation on public benchmarks.

I confirm that this manuscript is original, has not been published previously,
and is not under consideration for publication elsewhere. There are no conflicts
of interest to declare, and no external funding was received. All code, the
collected trajectory corpus, and the scripts that regenerate every reported
number and figure are released with the project repository to support full
reproducibility. As sole author, I am responsible for all aspects of the work.

Thank you for your time and consideration. I look forward to the reviewers'
feedback.

\vspace{0.8em}
\noindent Sincerely,

\vspace{0.4em}
\noindent Quang-Vinh Dang

\end{document}
